% Source: ME18_v1.html

\documentclass[11pt]{article}
\usepackage[margin=1in]{geometry}
\usepackage{amsmath}
\usepackage{amssymb}
\usepackage{siunitx}
\usepackage{enumitem}

\title{ME18: Heat and Calorimetry}
\author{}
\date{}

\begin{document}
\maketitle

\section*{ME18: Heat and Calorimetry}

\begin{enumerate}[leftmargin=*,label=\textbf{Question \arabic*.}]

\item \hfill \textit{10 pts}

Calculate the volume in liters of 0.8 moles of an ideal gas at STP (1 atm and 273 K). (1 L = 10$^{-3}$ m$^{3}$)



\emph{V} = \rule{2.5cm}{0.4pt} L  (Note the units!)

\textit{Correct answer:} 17.9214

\item \hfill \textit{10 pts}

An ideal gas has an initial pressure of \emph{p}$_{0}$ and volume of  1 m$^{3}$. It undergoes a process whereby the pressure changes to  (1.9)\emph{p}$_{0}$ while the Kelvin temperature increases by 85\%. Calculate the final volume in meters cubed.

\textit{Correct answer:} 0.9737

\item \hfill \textit{10 pts}

An ideal gas undergoes a process whereby the new pressure is 0.9 times the original pressure and the new volume is 0.7 times the original volume. If the original temperature was 297 K, calculate the new temperature in Kelvin.

\textit{Correct answer:} 187.11

\item \hfill \textit{10 pts}

A car tire has a gauge pressure of 211,466 Pa, volume of 3 L, and temperature of 5°C. After several hours of driving, the temperature of the air increases by 19C° and 9\% of the gas leaks out due to a nail stuck in the tread. Calculate the gauge pressure of the tire in psi at this time if the ambient atmospheric pressure is 1.00 × 10$^{5}$ Pa. Neglect any volume change of the tire. (1000 L = 1 m$^{3}$; 1 psi = 6895 Pa)  (Be careful with units!)



\emph{p}$_{gauge}$ = \rule{2.5cm}{0.4pt} psi

\textit{Correct answer:} 29.4134

\item \hfill \textit{10 pts}

Calculate the number of molecules in 77.1 attograms of C$_{4}$H$_{8}$OH. (1 attogram = 10$^{-18}$ g)

\textit{Correct answer:} 635,071.3217

\item \hfill \textit{10 pts}

Calculate the number of moles in 501 g of aluminum chloride, AlCl$_{3}$.

\textit{Correct answer:} 3.7573

\item \hfill \textit{10 pts}

The translational kinetic energy of 8 moles of an ideal gas is 64,356 J. Calculate the gas temperature in Kelvin.

\textit{Correct answer:} 645.0177

\item \hfill \textit{10 pts}

Calculate the total translational kinetic energy in J of 7.6 moles of an ideal gas at a temperature of 90°C.

\textit{Correct answer:} 34,407

\item \hfill \textit{10 pts}

Calculate the average speed (in m/s) of molecular oxygen at a temperature of 81.9 °C. (Molar mass of atomic oxygen is 16.0 g/mol.)

\textit{Correct answer:} 484.6685

\item \hfill \textit{10 pts}

Calculate the average value of the  $x$ -component of the \rule{2.5cm}{0.4pt} (in m/s) of molecular oxygen at a temperature of 69.4 °C. (Molar mass of atomic oxygen is 16.0 g/mol.)

\textit{Correct answer:} 0

\item \hfill \textit{10 pts}

Calculate the rms speed in m/s of molecular oxygen at a temperature of 39°C. (Molar mass of atomic oxygen is 16.0 g/mol.)

\textit{Correct answer:} 493.1522

\item \hfill \textit{10 pts}

Calculate the total internal energy in kJ of 2.4 mol of helium gas at a temperature of 155°C.



\emph{E}= \rule{2.5cm}{0.4pt} kJ

\textit{Correct answer:} 12.811

\item \hfill \textit{10 pts}

A gas has 2 degrees of freedom. By the equipartition of energy, calculate the internal energy change in kJ required to raise the temperature of 6.2 moles of the gas by 53 C°.

\textit{Correct answer:} 2.7321

\item \hfill \textit{10 pts}

Boltzman's constant \emph{k} is given by

    [ Select ]

  ["Avogadro's number", "the gas constant", "the number of degrees of freedom"]



    [ Select ]

  ["divided by", "multiplied by", "subtracted from"]



    [ Select ]

  ["Avogadro's number", "the gas constant", "the number of degrees of freedom", "the atomic mass unit"]

.

\begin{itemize}
  \item Avogadro's number
  \item the gas constant
  \item the number of degrees of freedom
  \item divided by
  \item multiplied by
  \item subtracted from
  \item Avogadro's number
  \item the gas constant
  \item the number of degrees of freedom
  \item the atomic mass unit
\end{itemize}

\item \hfill \textit{10 pts}

Tank A contains a monatomic gas, whereas tank B contains a diatomic gas, both systems having the same volume, pressure, number of moles, and temperature. If the same amount of energy Q is added to both systems, then the new temperature of gas A

    [ Select ]

  ["relative to gas B requires more information.", "is higher than that of gas B.", "is the same as that of gas B.", "is lower than that of gas B."]

 Further, the rms speed of the molecules of gas A

    [ Select ]

  ["is larger than that of gas B.", "is the same as that of gas B.", "is impossible to determine relative gas B without more information.", "is smaller than that of gas B."]

\begin{itemize}
  \item relative to gas B requires more information.
  \item is higher than that of gas B.
  \item is the same as that of gas B.
  \item is lower than that of gas B.
  \item is larger than that of gas B.
  \item is the same as that of gas B.
  \item is impossible to determine relative gas B without more information.
  \item is smaller than that of gas B.
\end{itemize}

\item \hfill \textit{10 pts}

There are three measures of speed in a gas. The highest speed is given by the

    [ Select ]

  ["most probable speed", "average speed", "root-mean-square speed", ""]

, whereas the lowest is given by the

    [ Select ]

  ["root-mean-square speed", "most probable speed", "average speed"]

.

\begin{itemize}
  \item most probable speed
  \item average speed
  \item root-mean-square speed
  \item root-mean-square speed
  \item most probable speed
  \item average speed
\end{itemize}

\end{enumerate}

\end{document}
